\paragraph{The published prime-two Euler-system input.}
For $p=2$ we use the rational, away-from-two case of
\cite[Theorem 0.8, pp.~317--318]{kato1999}, with the stronger
hypothesis in its condition 0.6(ii), to justify the application of
\cite[Theorem 13.4(2)]{kato}. For our non-CM elliptic representation
the image of $G_{\Q^{\rm ab}}$ contains an open subgroup of
$\mathrm{SL}_2(\Z_2)$, as stated in
\cite[Example 0.6.2, pp.~317--318]{kato1999}, using Serre's open-image
theorem. A quadratic character $\chi$ is trivial on $G_{\Q^{\rm ab}}$.
Thus the same subgroup acts on $V\chi$ and contains a matrix
$\left(\begin{smallmatrix}1&b\\0&1\end{smallmatrix}\right)$ with
$b\ne0$. Its fixed space over $\Q_2$ has dimension one. This supplies
the unipotent in $G_{\Q^{\rm ab}}$, not merely in the cyclotomic
subgroup. The latter subgroup also contains this open
$\mathrm{SL}_2$ subgroup, so $V\chi$ remains absolutely irreducible
over $\Q(\mu_{2^\infty})$. Indeed, a line invariant under both upper
and lower nonidentity unipotents cannot exist, even over an algebraic
closure. No inference from irreducibility over $G_\Q$ alone is needed.

We check the index set and the supplied class in that theorem
explicitly. Put $N_K=2Nd$ and $N'_K=cd_s$. These integers are
coprime by (K5), since $d\mid F$. Section~0.1 of \cite{kato1999}
requires every index $m$ with $N_K\mid m$ and $(m,N'_K)=1$.
The enlarged definition (K8) supplies all these classes. They
originate in the integral $j_*$ group over
$\mathbf Z[\zeta_m,1/2]$ and map to full cohomology over
$\mathbf Z[\zeta_m,1/N_K]$. The varying construction support
$S_{Dm}$ is not an instruction to invert all its primes in this
last scheme. Apply the fixed lattice map $2^{c_E}$ from (K4) if
needed, identically at every level.

For a prime $\ell\nmid N'_K$, if $\ell\mid m$ then
$S_{Dm\ell}=S_{Dm}$, so Proposition~8.12 gives norm factor one.
If $\ell\nmid m$, the condition $N_K\mid m$ implies
$\ell\nmid2Nd$, and hence
$S_{Dm\ell}\setminus S_{Dm}=\{\ell\}$. The equivariance identity
before (K9) gives the factor
$P_\ell(\ell^{-1}\sigma_\ell^{-1})=E_\ell(\chi)$, with
$P_\ell(X)=1-a_\ell\chi(\ell)X+\ell X^2$. These are exactly the
two cases of (0.1.1) in \cite{kato1999}. Smoothing primes are
excluded by $N'_K$ and need not be inverted.

For the class in Section~0.7 of that source take its parameter
$d=1$; this parameter is distinct from our odd conductor part $d$.
For every $n\geq1$,
\[
 S_{DN_K2^n}=S_{D2^n}=\operatorname{prime}(2Nd).
\]
Proposition~8.12 therefore gives
\[
 \operatorname{cor}_{\mathbf Q(\zeta_{DN_K2^n})/
                          \mathbf Q(\zeta_{D2^n})}
             z_{DN_K2^n}=z_{D2^n}.
\]
Transitivity of corestriction, through the two intermediate fields
in the common upper field, now yields
\[
 \operatorname{cor}_{\mathbf Q(\zeta_{N_K2^n})/
                          \mathbf Q(\zeta_{2^n})}y_{N_K2^n}
 =\operatorname{cor}_{\mathbf Q(\zeta_{D2^n})/
                          \mathbf Q(\zeta_{2^n})}
                         (z_{D2^n}\otimes e_\chi)
 =y_{2^n}.
\]
No Euler factor occurs because the supports are equal. No field
degree occurs because this is corestriction of the actual upper
class, not corestriction of a restriction. These equalities commute
with transition maps, and $n\geq1$ is a cofinal tail. Write
$\xi_K$ for the Section~0.7 class, to distinguish it from the
modular-symbol combination $\xi$. It is exactly the full-cohomology
image of $y_{2^\infty}$, multiplied by $2^{c_E}$ if the fixed
lattice map was needed. That factor is a unit at every height-one
prime under consideration, so its coordinate ideal is precisely
$J_{\mathfrak p}$. The preceding
nonvanishing argument supplies a nontorsion class on each sign
component, purity gives $w=-1$, and oddness gives condition
0.6(iv$_{\mathfrak p}$). At the height-one primes used here, $2$ is
invertible and the localized Iwasawa algebra is a DVR. Theorem 0.8
then bounds the coordinate ideal of this actual retained class. Its
strong residual and integral-cokernel assumptions concern primes
containing $2$, which are not used here.

